\chapter{Red Blood Cell Indices}

\section*{What Is This Test?}
Red blood cell indices are calculated parameters derived from the measured hemoglobin concentration, red blood cell count, and hematocrit that characterize the average physical properties of an individual's red blood cells. The four standard indices are: mean corpuscular volume (MCV), which describes average erythrocyte size; mean corpuscular hemoglobin (MCH), the average hemoglobin content per red blood cell; mean corpuscular hemoglobin concentration (MCHC), the average hemoglobin concentration within a given volume of packed red cells; and red cell distribution width (RDW), which quantifies the variation in red blood cell size (anisocytosis). Together, these indices form the cornerstone of anemia classification and differential diagnosis. By integrating the RBC indices, clinicians can determine whether an anemia is microcytic (MCV <80 fL), normocytic (MCV 80--100 fL), or macrocytic (MCV >100 fL), and whether it is hypochromic (low MCH and MCHC) or normochromic, thereby narrowing the diagnostic possibilities and directing appropriate further testing \cite{tierney2023current}.

\section*{Alternative Names}
\begin{itemize}
  \item RBC Indices
  \item Erythrocyte Indices
  \item Red Cell Indices
  \item MCV, MCH, MCHC, RDW
  \item Wintrobe Indices
  \item Corpuscular Indices
\end{itemize}
\section*{Principle}
Red blood cell indices are calculated by the automated hematology analyzer from directly measured parameters. MCV is derived from the volume distribution histogram generated by impedance counting: the mean of the erythrocyte volume distribution is the MCV. Some optical analyzers measure MCV directly by laser light scatter, correlating the forward scatter angle with cell volume. MCH is calculated as: MCH (pg) = Hemoglobin (g/dL)  $\times$ 10 /  RBC count  ($\times$ 10\textsuperscript{6}/$\mu$L). MCHC is calculated as: MCHC (g/dL) = Hemoglobin (g/dL)  $\times$ 100 /  Hematocrit (\%) . \cite{tietz2012textbook} RDW is a statistical parameter computed from the RBC volume distribution histogram. RDW-CV (coefficient of variation) is the standard deviation of erythrocyte volumes divided by the MCV, expressed as a percentage. RDW-SD is the standard deviation of the volume distribution in fL. Beckman Coulter, Sysmex, and Siemens analyzers all provide these calculated indices as standard CBC parameters. Quality control is maintained by daily calibration with reference control materials and proficiency testing through external quality assurance programs.

\section*{History}
The concept of red blood cell indices was pioneered by Maxwell Wintrobe in the 1930s. In a landmark 1934 paper, Wintrobe described the MCV, MCH, and MCHC as objective, quantitative parameters for classifying anemias, a revolutionary advance over the qualitative descriptions of ``microcytic'' and ``macrocytic'' anemic blood smears \cite{wintrobe1934anemia}. Wintrobe manually measured hematocrit by centrifugation (the Wintrobe tube method), RBC count by hemocytometer, and hemoglobin by colorimetric methods, then computed the indices using his formulas. These ``Wintrobe indices'' became the standard for anemia classification for decades. The RDW was introduced later, in the 1980s, following the development of automated analyzers that generated volume histograms. Currently, the indices are calculated automatically by every modern hematology analyzer, and they remain essential to hematology practice, as originally envisioned by Wintrobe nearly a century ago \cite{brugnara2013red}.

\section*{Why This Test Is Ordered}
\begin{itemize}
  \item Classification of anemia: MCV categorizes anemia as microcytic (<80 fL), normocytic (80--100 fL), or macrocytic (>100 fL), which directs the differential diagnosis \cite{kaushansky2021williams}
  \item Differentiation of iron deficiency anemia from thalassemia trait: both are microcytic, but iron deficiency shows elevated RDW while thalassemia trait typically shows normal RDW
  \item Screening for iron deficiency: low MCV and MCH precede the development of overt anemia and may be the earliest hematologic indicators of iron-deficient erythropoiesis
  \item Evaluation of macrocytosis: elevated MCV raises suspicion for vitamin B12 or folate deficiency, chronic liver disease, alcohol abuse, hypothyroidism, or myelodysplastic syndrome
  \item Detection of hereditary spherocytosis: elevated MCHC (>36 g/dL) with spherocytes on peripheral smear is characteristic
  \item Diagnosis of cold agglutinin disease: spuriously elevated MCV and MCHC with falsely low RBC count due to red cell clumping
  \item Monitoring response to hematinic therapy: normalization of MCV and RDW as iron, B12, or folate stores are replenished
  \item Prognostic assessment: elevated RDW has been independently associated with increased mortality in cardiovascular disease, heart failure, and critical illness
\end{itemize}

\section*{Primary vs Secondary Test}
Red blood cell indices are primary parameters calculated from the CBC and are essential for the classification and differential diagnosis of anemia.

\section*{How This Test Is Performed}
The red blood cell indices are calculated automatically from a single venous blood sample collected in a lavender-top (EDTA) tube. No additional sample or separate testing is required beyond the standard CBC. The blood sample is processed on an automated hematology analyzer, which measures hemoglobin concentration spectrophotometrically, counts red blood cells by impedance or optical methods, and determines the cell volume distribution. The indices are then computed by the analyzer's software and reported as part of the CBC panel \cite{bain2017dacie}. The venipuncture procedure takes less than 5 minutes, and automated analysis is completed within 1--2 minutes.

\section*{What Preparation Is Needed}
No special preparation is required. The red blood cell indices are derived from the same blood sample used for the CBC, so no fasting or other preparation is necessary. Patients should avoid vigorous exercise immediately before testing. Medications, particularly iron supplements, vitamin B12, folate, erythropoietin, recent blood transfusion, and chemotherapy should be disclosed to the healthcare provider, as they can affect the indices.

\section*{Contraindications and Precautions}
\begin{itemize}
  \item No absolute contraindications. Several pre-analytical and analytical factors can produce spurious red cell index results: cold agglutinins cause red cell clumping, resulting in a falsely low RBC count, falsely high MCV, and falsely elevated MCHC; severe hyperglycemia (>600 mg/dL) can cause osmotic swelling of red blood cells in the sample, increasing MCV; delayed sample processing (>24 hours) causes red cell swelling and increased MCV; marked leukocytosis (>100,000/$\mu$L) can interfere with some analyzers' RBC/MCV measurements; severe lipemia may interfere with hemoglobin measurement and thus affect MCH and MCHC; hemolyzed samples give unreliable RBC counts and invalid indices. EDTA is the standard anticoagulant. Samples should be analyzed within 4 hours at room temperature or 24 hours if refrigerated at 2--8\textdegree{}C.
\end{itemize}
\section*{Risks}
The risks are those of routine venipuncture: mild pain or stinging at the needle site, bruising, hematoma, and rarely, vasovagal syncope or infection. No additional risks apply specifically to the indices, as they are calculated from the CBC sample.

\section*{What Happens After the Test}
Results are available within 30--60 minutes as part of the CBC. The healthcare provider interprets the indices in combination with hemoglobin, hematocrit, and the peripheral blood smear. The pattern of MCV and RDW directs the initial diagnostic approach: microcytic anemia with high RDW suggests iron deficiency (confirmed by low ferritin); microcytic anemia with normal RDW suggests thalassemia trait (confirmed by hemoglobin electrophoresis); macrocytic anemia suggests B12 or folate deficiency (confirmed by serum levels); normocytic anemia with high RDW suggests early iron deficiency, mixed deficiency, or hemoglobinopathy \cite{wallach2022interpretation}. Serial monitoring of MCV and RDW tracks response to iron, B12, or folate therapy.

\section*{Understanding Results}

\subsection*{Below Normal (MCV <80 fL --- Microcytosis)}
\begin{itemize}
  \item \textbf{Iron deficiency anemia:} Low MCV and MCH with elevated RDW, low serum ferritin (<30 ng/mL), low transferrin saturation (<16\%), and elevated TIBC.
  \item \textbf{Thalassemia trait (alpha or beta):} Low MCV and MCH with normal or minimally elevated RDW, normal or elevated RBC count, normal iron studies, and abnormal hemoglobin electrophoresis (elevated HbA2 in beta-thal trait).
  \item \textbf{Anemia of chronic disease:} Mildly low MCV (usually 75--80 fL) with low serum iron, low TIBC, normal or elevated ferritin, and inflammatory markers.
  \item \textbf{Sideroblastic anemia:} Dimorphic red cell population with low MCV and elevated RDW, ring sideroblasts on bone marrow iron stain, elevated serum iron and ferritin.
  \item \textbf{Lead poisoning:} Microcytic hypochromic anemia with basophilic stippling on peripheral smear, elevated blood lead level (>5 $\mu$g/dL), and elevated zinc protoporphyrin.
  \item \textbf{Hemoglobin E disease and hemoglobin H disease:} Microcytic anemia with normal ferritin and abnormal hemoglobin patterns on electrophoresis.
\end{itemize}

\subsection*{Above Normal (MCV >100 fL --- Macrocytosis)}
\begin{itemize}
  \item \textbf{Vitamin B12 (cobalamin) deficiency:} Marked macrocytosis (MCV often 110--130 fL) with elevated RDW, hypersegmented neutrophils (>5\% of neutrophils with 5 or more lobes), low serum B12 (<200 pg/mL), and elevated methylmalonic acid and homocysteine.
  \item \textbf{Folate deficiency:} Macrocytosis indistinguishable from B12 deficiency by routine indices, with low serum folate (<2 ng/mL) or low RBC folate, and elevated homocysteine but normal methylmalonic acid.
  \item \textbf{Chronic liver disease:} Mild macrocytosis (MCV 100--110 fL) without characteristic B12/folate laboratory abnormalities; target cells may be present on smear.
  \item \textbf{Chronic alcohol abuse:} Direct toxic effect on erythropoiesis plus associated folate deficiency; MCV may remain elevated for weeks after abstinence.
  \item \textbf{Hypothyroidism:} Mild macrocytosis that resolves with thyroid hormone replacement.
  \item \textbf{Myelodysplastic syndrome:} Macrocytosis (often MCV 100--115 fL) with elevated RDW, cytopenias in other cell lines, and dysplastic morphology on bone marrow examination.
  \item \textbf{Drug-induced macrocytosis:} Hydroxyurea, zidovudine (AZT), methotrexate, phenytoin, trimethoprim, and certain chemotherapeutic agents.
  \item \textbf{Reticulocytosis:} Large numbers of young, larger red blood cells can shift the MCV upward; the reticulocyte count should be checked.
\end{itemize}

\subsection*{MCHC Above Normal}
\begin{itemize}
  \item \textbf{Hereditary spherocytosis:} MCHC >36 g/dL, spherocytes on peripheral smear, increased osmotic fragility, negative direct Coombs test.
  \item \textbf{Cold agglutinin disease:} Spurious elevation of MCHC with red cell clumping on smear; warming sample to 37\textdegree{}C resolves the abnormality.
\end{itemize}

\subsection*{RDW Above Normal}
\begin{itemize}
  \item \textbf{Iron deficiency anemia:} RDW is typically elevated early, before anemia develops, and is a sensitive marker.
  \item \textbf{Mixed deficiency states:} Combined iron and B12/folate deficiency produces a dimorphic blood picture with markedly elevated RDW.
  \item \textbf{Post-transfusion or post-treatment:} Dimorphic populations of native and transfused or newly produced cells.
  \item \textbf{Myelofibrosis:} Teardrop cells and other poikilocytes contributing to anisocytosis.
\end{itemize}

\section*{Normal Results}
\begin{itemize}
  \item \textbf{MCV (Mean Corpuscular Volume):} Adults: 80--100 fL; Children (1--6 years): 72--87 fL; Newborns: 98--122 fL (first week) \cite{fischbach2023manual}
  \item \textbf{MCH (Mean Corpuscular Hemoglobin):} Adults: 27--33 pg; Children (1--6 years): 23--31 pg; Newborns: 33--39 pg
  \item \textbf{MCHC (Mean Corpuscular Hemoglobin Concentration):} Adults: 32--36 g/dL; Children: 31--36 g/dL; Newborns: 30--36 g/dL
  \item \textbf{RDW-CV (Red Cell Distribution Width):} Adults: 11.5--14.5\%
  \item \textbf{RDW-SD:} Adults: 38--47 fL (platform-dependent)
\end{itemize}

\section*{Follow-up Tests}
\begin{itemize}
  \item \textbf{Peripheral blood smear:} Direct morphologic examination to confirm index findings --- microcytosis, macrocytosis, hypochromia, anisocytosis, target cells, spherocytes
  \item \textbf{Iron studies:} Serum iron, ferritin, total iron-binding capacity, transferrin saturation --- to confirm iron deficiency and distinguish from anemia of chronic disease
  \item \textbf{Vitamin B12 and serum folate:} To diagnose megaloblastic anemia when MCV is elevated
  \item \textbf{Methylmalonic acid and homocysteine:} More sensitive markers for B12 and folate deficiency when serum levels are borderline
  \item \textbf{Hemoglobin electrophoresis:} To detect thalassemia trait (elevated HbA2) and hemoglobinopathies (HbS, HbC, HbE)
  \item \textbf{Direct antiglobulin test (Coombs test):} When spherocytes are seen on smear with elevated MCHC
  \item \textbf{Osmotic fragility test:} For suspected hereditary spherocytosis
  \item \textbf{Serum erythropoietin level:} When macrocytosis without megaloblastic features is present
  \item \textbf{Thyroid function tests (TSH, free T4):} When macrocytosis is unexplained
  \item \textbf{Liver function tests:} When liver disease is suspected as cause of macrocytosis
  \item \textbf{Bone marrow aspiration and biopsy:} When myelodysplastic syndrome or other marrow disorder is suspected
  \item \textbf{Blood lead level and zinc protoporphyrin:} If lead poisoning is suspected with microcytosis and basophilic stippling
  \item \textbf{Reticulocyte count:} To rule out reticulocytosis as contributor to macrocytosis
\end{itemize}

\section*{References}
\bibliographystyle{unsrtnat}
\bibliography{references}

\clearpage
\section*{Regulatory Status and Authorization Date}
This chapter describes a test category, not a specific commercial product. FDA, CDSCO, EU IVDR, and MHRA approval or registration dates therefore cannot be assigned without the assay manufacturer, product name, instrument, and intended use. For laboratory-developed tests, an individual agency approval date may not exist. See the Preface for the interpretation of regulatory status.

\clearpage
